\begin{abstract}
This paper proposes a framework for separating shock-driven and exhaustion-driven termination mechanisms in armed conflict. We introduce two complementary composite indices---the Decisive Shock Score (DSS) and the Strategic Exhaustion Score (SES)---defined as transparent weighted averages of normalized, dimensionless component scores, and apply them to a dataset of 4,812 wars. Many DSS components require post-hoc information; we therefore distinguish observed DSS (explanatory) from a structural DSS proxy (partially predictive, using only pre-war features). Among the 91 wars with sufficient battle-level data, 2.2\% are classified as decisive shock and 22.0\% as strategic exhaustion; the remainder are uncertain or mixed. A dynamical simulation model demonstrates that shock and exhaustion mechanisms can coexist in a bounded state-space model: when calibrated to seven historical cases, the mechanism classifier agrees with our pre-specified case-study labels in six cases (86\%). We treat this as an internal consistency check, not an independent validation result. A random forest classifier achieves 72.7\% accuracy predicting war duration from material-capability features, suggesting nonlinear structure in these predictors. Independent blind validation of the mechanism classifier remains future work.
\end{abstract}
